\section{Related Work: A Cost Shape, Not a Speedup Number}
\label{sec:related}

\paragraph{Diffusion language models and their speedup claims}
LLaDA-8B~\cite{llada2025} and DiffusionGemma-26B~\cite{diffusiongemma} are the
two diffusion decoders whose serving cost this paper measures. Both are
presented as low-latency routes to structured generation, and the speedup
ratios that accompany them are measured against autoregressive baselines
assembled by the same authors. Our motivating observation
(\S\ref{sec:intro}) is not that these claims are wrong but that the baseline
is not a neutral object: the decoding configuration, serving runtime, parser,
and metric each admit a multiple, and we found five defects in our
\emph{own} harness --- four of them biasing the same arm --- before any
architectural conclusion could hold. Our delta is to put both arms on one
pinned stack (vLLM \vllmVersion), fix those defects, and report a fitted cost
\emph{shape} rather than a single speedup number. DiffusionGemma-26B-A4B is
the diffusion arm; LLaDA-8B under a second, diffusion-native runtime is a
control whose cache bug confines it to an inconclusive result
(\S\ref{sec:threats}).

\paragraph{Constrained and structured decoding}
The standard answer to fixed-schema generation is grammar- or
regex-constrained sampling: outlines~\cite{outlines2024},
xgrammar~\cite{xgrammar2024}, llguidance~\cite{llguidance}, and
jsonformer~\cite{jsonformer} each force every sampled token into the language
of the target schema. vLLM \vllmVersion~\cite{vllm2023} ships xgrammar and
llguidance on its structured-output path. Our task --- 12-field JSON
extraction from LOGHUB log lines~\cite{loghub} (BGL and HDFS) --- is exactly
the workload these tools target, and we do \emph{not} include a
constrained-decoding baseline. That is a stated limitation (issue \#72):
constrained decoding does not change the autoregressive cost \emph{shape} ---
it is still a line in output length with no floor --- but it can raise schema
validity and thereby tighten the quality-matched pair, which would move the
crossover. We name the gap rather than close it.

\paragraph{Speculative decoding}
Medusa~\cite{medusa2024}, EAGLE~\cite{eagle2024}, and lookahead
decoding~\cite{lookahead2024} attack the same per-token slope
(\arSlope~ms/token) that this paper's argument turns on, by drafting and
verifying multiple tokens per forward pass. Our delta is that speculative
decoding changes the \emph{slope} of the autoregressive line, not its
\emph{shape}: the cost remains linear in output length with no floor, so the
crossover with the diffusion staircase moves rightward rather than
vanishing. We do not measure speculative decoding on the pinned stack; as
with constrained decoding, it is a baseline the fair comparison
(\S\ref{sec:fair}) would need before any deployment claim is final.

\paragraph{Benchmarking methodology and reproducibility}
Prior work on measurement discipline in AI systems~\cite{hoefler2021}
argues that benchmarking infrastructure carries biases invisible from the
reported numbers alone. Our contribution is narrower and more concrete: the
five harness defects in \S\ref{sec:setup} were each found by re-reading
committed artifacts --- not by running new experiments --- and four of the
five biased against the same arm. The lesson is that an experimenter who
prefers one architecture builds a harness that fails asymmetrically, and the
remedy is an audit checklist (\artifact{CHECKLIST.md}) whose items each take
under five minutes and each caught a defect that had already moved a
headline number.
